\chapter{Gaussian integrals and Hermite polynomials}
\label{ch:gaussian-hermite}

The paper's first theorem says that the Ornstein--Uhlenbeck channel multiplies the
Hermite moment of degree $m$ by $e^{-mt}$, exactly, for every law with finite moments
of that order. Its proof is four lines long once three classical facts are in place:
the Gaussian generating function, the generating function of the Hermite polynomials,
and Mehler's formula. This chapter builds those facts from the Gaussian integral
upwards, proves Mehler's formula and the paper's Theorem 1 in full, derives the two
exact consequences the paper states after it, works out the Wick (Isserlis) rule and
the second-moment lemma that produce the constants $\tfrac16$ and $\tfrac1{24}$,
explains when a density ratio belongs to $L^2(\gamma)$ and why the paper does not need
it to, and defines the first Edgeworth correction that the chapter on
misspecification uses.

\section{Gaussian integrals}

\begin{lemma}[The scalar Gaussian integral and its moments]
\label{lem:gauss-scalar}
$\int_{\Reals}e^{-x^2/2}\dd x=\sqrt{2\pi}$. Writing $\gamma(x)=(2\pi)^{-1/2}e^{-x^2/2}$
for the \term{standard Gaussian density} and $Z\sim\gamma$,
\begin{equation}
  \E[e^{sZ}]=e^{s^2/2},\qquad \E[Z^{2k}]=(2k-1)!!=\frac{(2k)!}{2^kk!},\qquad \E[Z^{2k+1}]=0.
  \label{eq:gauss-moments}
\end{equation}
\end{lemma}
\begin{proof}
Squaring the integral and passing to polar coordinates,
\[
  \Bigl(\int e^{-x^2/2}\dd x\Bigr)^2=\int_0^{2\pi}\!\!\int_0^\infty e^{-r^2/2}r\,\dd r\,\dd\phi=2\pi .
\]
Completing the square, $sz-z^2/2=-(z-s)^2/2+s^2/2$, so
\[
  \E[e^{sZ}]=e^{s^2/2}(2\pi)^{-1/2}\int e^{-(z-s)^2/2}\dd z=e^{s^2/2}.
\]
Expanding both sides in $s$, $\sum_n\E[Z^n]s^n/n!=\sum_k s^{2k}/(2^kk!)$, and comparing
coefficients of $s^{2k}$ gives $\E[Z^{2k}]=(2k)!/(2^kk!)$, which equals
$(2k-1)!!=1\cdot3\cdots(2k-1)$; odd coefficients vanish. The exchange of sum and expectation is justified by
$\sum_n|s|^n|Z|^n/n!=e^{|s||Z|}$, which is integrable.
\end{proof}

So $\E Z^2=1$, $\E Z^4=3$, $\E Z^6=15$. The number $(2k-1)!!$ is the number of ways of
pairing $2k$ objects, which is the content of the next section.

\begin{lemma}[The multivariate Gaussian integral]
\label{lem:gauss-multi}
For a symmetric positive definite $\mat A$ on $\Reals^{\nsites}$ and $\vect b\in\Reals^{\nsites}$,
\begin{equation}
  \int_{\Reals^{\nsites}}\exp\bigl(-\tfrac12\vect x^{\!\top}\mat A\vect x+\vect b^{\!\top}\vect x\bigr)\dd\vect x
  =(2\pi)^{\nsites/2}(\det\mat A)^{-1/2}\exp\bigl(\tfrac12\vect b^{\!\top}\mat A^{-1}\vect b\bigr).
  \label{eq:gauss-multi}
\end{equation}
Consequently $\normal{\vect m}{\Sig}$ has density
$(2\pi)^{-\nsites/2}(\det\Sig)^{-1/2}\exp(-\tfrac12(\vect x-\vect m)^{\!\top}\Sig^{-1}(\vect x-\vect m))$
and generating function $\E[e^{\vect s^{\!\top}\vect X}]=\exp(\vect s^{\!\top}\vect m+\tfrac12\vect s^{\!\top}\Sig\vect s)$.
\end{lemma}
\begin{proof}
Complete the square,
\[
  -\tfrac12\vect x^{\!\top}\mat A\vect x+\vect b^{\!\top}\vect x
  =-\tfrac12(\vect x-\mat A^{-1}\vect b)^{\!\top}\mat A(\vect x-\mat A^{-1}\vect b)+\tfrac12\vect b^{\!\top}\mat A^{-1}\vect b,
\]
shift, and diagonalise $\mat A=\mat O\diag(\lambda_k)\mat O^{\!\top}$ with $\mat O$ orthogonal;
in the rotated coordinates $\vect y=\mat O^{\!\top}\vect x$ the integral factorises into
\[
  \prod_k\int e^{-\lambda_ky_k^2/2}\dd y_k=\prod_k\sqrt{2\pi/\lambda_k}=(2\pi)^{\nsites/2}(\det\mat A)^{-1/2}.
\]
The density and the generating function follow with $\mat A=\Sig^{-1}$.
\end{proof}

The standard Gaussian on $\Reals^{\nsites}$, written $\gamma$ throughout, is the case
$\Sig=\Id$; its coordinates are independent standard Gaussians, which is the fact
that reduces every multivariate Hermite statement below to the scalar one.

\section{Isserlis' theorem (Wick's rule)}

\begin{theorem}[Isserlis]
\label{thm:isserlis}
Let $\vect X\sim\normal{\vect0}{\Sig}$. For any indices $i_1,\dots,i_{2k}$ (repetitions
allowed),
\begin{equation}
  \E[X_{i_1}\cdots X_{i_{2k}}]=\sum_{\text{pairings }\pi}\ \prod_{\{a,b\}\in\pi}\Sigma_{i_ai_b},
  \label{eq:isserlis}
\end{equation}
the sum over the $(2k-1)!!$ ways of partitioning $\{1,\dots,2k\}$ into $k$ unordered
pairs; products of an odd number of coordinates have expectation zero.
\end{theorem}
\begin{proof}[Proof sketch]
By \cref{lem:gauss-multi}, $\E[e^{\vect s^{\!\top}\vect X}]=\exp(\tfrac12\sum_{ab}\Sigma_{ab}s_as_b)$.
The moment $\E[X_{i_1}\cdots X_{i_{2k}}]$ is the mixed derivative
$\partial_{s_{i_1}}\cdots\partial_{s_{i_{2k}}}$ at $\vect s=\vect 0$ of the left side.
On the right side, expand the exponential: the term of order $k$ is
$\tfrac1{k!}(\tfrac12\sum_{ab}\Sigma_{ab}s_as_b)^k$, a homogeneous polynomial of degree
$2k$ in $\vect s$; differentiating it $2k$ times and setting $\vect s=\vect 0$ picks out,
for each way of assigning the $2k$ differentiations to the $k$ factors in pairs, the
product of the corresponding $\Sigma$ entries. Each unordered pairing arises
$k!\,2^k$ times (order of the factors, order within each pair), which cancels the
$\tfrac1{k!}\tfrac1{2^k}$ in front. Odd derivatives vanish because the right side is
even in $\vect s$.
\end{proof}

For $k=2$ this is $\E[X_iX_jX_kX_l]=\Sigma_{ij}\Sigma_{kl}+\Sigma_{ik}\Sigma_{jl}+\Sigma_{il}\Sigma_{jk}$,
the three pairings that appeared in the fourth-cumulant formula of the chapter on
moments and cumulants: for a Gaussian vector the fourth moment is exactly the sum over
pairings, so the fourth cumulant, which subtracts that sum, vanishes.

\section{Probabilists' Hermite polynomials}

\begin{definition}[Rodrigues formula]
The \term{probabilists' Hermite polynomial} of degree $n$ is
\begin{equation}
  \Hen{n}(x)=(-1)^ne^{x^2/2}\frac{\dd^n}{\dd x^n}e^{-x^2/2},\qquad n=0,1,2,\dots
  \label{eq:rodrigues}
\end{equation}
\end{definition}

Differentiating $e^{-x^2/2}$ once gives $-xe^{-x^2/2}$, twice $(x^2-1)e^{-x^2/2}$, and
so on; each derivative is a polynomial times $e^{-x^2/2}$, and the sign $(-1)^n$ makes
the polynomial monic. The first seven are
\begin{equation}
\begin{aligned}
  \Hen{0}(x)&=1, & \Hen{1}(x)&=x, & \Hen{2}(x)&=x^2-1,\\
  \Hen{3}(x)&=x^3-3x, & \Hen{4}(x)&=x^4-6x^2+3, & \Hen{5}(x)&=x^5-10x^3+15x,\\
  \Hen{6}(x)&=x^6-15x^4+45x^2-15.
\end{aligned}
\label{eq:he-list}
\end{equation}
$\Hen{n}$ has the parity of $n$. The coefficient pattern is the Gaussian moment
sequence with alternating signs: $\Hen{n}(x)=\sum_k(-1)^k\binom{n}{2k}(2k-1)!!\,x^{n-2k}$,
which the generating function below explains.

\begin{proposition}[Generating function]
\label{prop:he-gf}
For all real $x$ and $s$,
\begin{equation}
  e^{sx-s^2/2}=\sum_{n\ge0}\Hen{n}(x)\frac{s^n}{n!}.
  \label{eq:he-gf}
\end{equation}
\end{proposition}
\begin{proof}
Write $e^{-(x-s)^2/2}=e^{-x^2/2}e^{sx-s^2/2}$. The left side, as a function of $s$,
is entire, and its Taylor expansion at $s=0$ is
$\sum_n\frac{s^n}{n!}\frac{\dd^n}{\dd s^n}e^{-(x-s)^2/2}\big|_{s=0}$. Since
$\frac{\dd}{\dd s}f(x-s)=-f'(x-s)$, the $n$-th derivative at $s=0$ is
$(-1)^n\frac{\dd^n}{\dd x^n}e^{-x^2/2}=\Hen{n}(x)e^{-x^2/2}$ by \cref{eq:rodrigues}.
Dividing by $e^{-x^2/2}$ gives \cref{eq:he-gf}.
\end{proof}

Two identities follow by differentiating \cref{eq:he-gf}. Differentiating in $x$,
$s\,e^{sx-s^2/2}=\sum_n\Hen{n}'(x)s^n/n!$, and comparing coefficients,
\begin{equation}
  \Hen{n}'(x)=n\Hen{n-1}(x).
  \label{eq:he-derivative}
\end{equation}
Differentiating in $s$, $(x-s)e^{sx-s^2/2}=\sum_n\Hen{n}(x)s^{n-1}/(n-1)!$, and
comparing coefficients of $s^n/n!$,
\begin{equation}
  \Hen{n+1}(x)=x\Hen{n}(x)-n\Hen{n-1}(x),
  \label{eq:he-recursion}
\end{equation}
the three-term recursion that \code{src/markov\_score/theory.py} (function \code{he})
uses to evaluate the polynomials. Expanding $e^{sx}e^{-s^2/2}$ as a product of two
series and collecting $s^n$ gives the explicit coefficient formula stated above.

\begin{proposition}[Orthogonality]
\label{prop:he-orth}
$\E_\gamma[\Hen{n}(Z)\Hen{m}(Z)]=n!\,\delta_{nm}$.
\end{proposition}
\begin{proof}
Multiply two copies of \cref{eq:he-gf} with parameters $s$ and $u$ and take the
expectation under $\gamma$:
\[
  \E_\gamma\bigl[e^{sZ-s^2/2}e^{uZ-u^2/2}\bigr]=e^{-(s^2+u^2)/2}\E_\gamma[e^{(s+u)Z}]
  =e^{-(s^2+u^2)/2}e^{(s+u)^2/2}=e^{su}=\sum_{n\ge0}\frac{s^nu^n}{n!}.
\]
On the other side the product of the two series is
$\sum_{n,m}\E_\gamma[\Hen{n}\Hen{m}]\,s^nu^m/(n!\,m!)$. Comparing the coefficient of
$s^nu^m$ gives $\E_\gamma[\Hen{n}\Hen{m}]/(n!\,m!)=\delta_{nm}/n!$, which is the claim. The exchange of expectation and double sum is justified as in
\cref{lem:gauss-scalar}, the double series being dominated by $e^{(|s|+|u|)|Z|}$.
\end{proof}

So $\{\Hen{n}/\sqrt{n!}\}_{n\ge0}$ is an orthonormal system in $L^2(\gamma)$, the
space of functions $f$ with $\E_\gamma[f(Z)^2]<\infty$. It is also complete: the
polynomials are dense in $L^2(\gamma)$ (one proof: if $f\in L^2(\gamma)$ is orthogonal to
every polynomial then the entire function $s\mapsto\E_\gamma[f(Z)e^{sZ}]$ has every
derivative equal to zero at the origin, hence the Fourier transform of $f\gamma$ vanishes,
hence $f=0$). A function $f\in L^2(\gamma)$ therefore has the \term{Hermite expansion}
\begin{equation}
  f=\sum_{n\ge0}\frac{\E_\gamma[f\Hen{n}]}{n!}\Hen{n},\qquad
  \E_\gamma[f^2]=\sum_{n\ge0}\frac{\E_\gamma[f\Hen{n}]^2}{n!},
  \label{eq:hermite-expansion}
\end{equation}
with convergence in $L^2(\gamma)$ and the Parseval identity on the right. The
\term{Hermite moment} of degree $n$ of a law $P$ is $\E_P[\Hen{n}]$; when $P$ has a
density and $P/\gamma\in L^2(\gamma)$, these are the expansion coefficients of the
density ratio: $P/\gamma=\sum_n\E_P[\Hen{n}]\Hen{n}/n!$.

\section{Products of Hermite polynomials and Wick products}

\begin{definition}[Multivariate Hermite products]
For a multi-index $\nu\in\mathbb N^{\nsites}$,
\[
  H_\nu(\vx)=\prod_{i=1}^{\nsites}\Hen{\nu_i}(x_i),\qquad |\nu|=\sum_i\nu_i,\qquad \nu!=\prod_i\nu_i!,
\]
and $|\nu|$ is the \term{degree} of $H_\nu$.
\end{definition}

Under $\gamma$ on $\Reals^{\nsites}$ the coordinates are independent, so the expectation
of a product factorises and \cref{prop:he-orth} gives at once
\begin{equation}
  \E_\gamma[H_\nu H_\mu]=\prod_i\E_\gamma[\Hen{\nu_i}\Hen{\mu_i}]=\prod_i\nu_i!\,\delta_{\nu_i\mu_i}=\nu!\,\delta_{\nu\mu}.
  \label{eq:multi-orth}
\end{equation}
This is the normalisation the paper states in Section 3; $\{H_\nu/\sqrt{\nu!}\}$ is an
orthonormal basis of $L^2(\gamma)$ on $\Reals^{\nsites}$ (completeness by the same
argument coordinatewise), and the \term{chaos expansion} of $f\in L^2(\gamma)$ is
$f=\sum_\nu\E_\gamma[fH_\nu]H_\nu/\nu!$. The paper writes $H_\nu$ with a capital $H$;
it is the product of probabilists' polynomials, not a physicists' polynomial (see the
chapter on conventions).

The same objects appear in the paper's proofs in a second notation, the
\term{Wick product} ${:}x_{i_1}\cdots x_{i_n}{:}$, defined by the multivariate
generating function
\begin{equation}
  e^{\vect s^{\!\top}\vx-|\vect s|^2/2}=\sum_{n\ge0}\frac1{n!}\sum_{i_1,\dots,i_n}s_{i_1}\cdots s_{i_n}\,{:}x_{i_1}\cdots x_{i_n}{:}\,.
  \label{eq:wick-gf}
\end{equation}
Since the left side is $\prod_ie^{s_ix_i-s_i^2/2}=\prod_i\sum_k\Hen{k}(x_i)s_i^k/k!$,
comparing the coefficient of a monomial $\prod_is_i^{\nu_i}$ (which appears
$n!/\nu!$ times among the ordered tuples $(i_1,\dots,i_n)$) shows
\begin{equation}
  {:}x_{i_1}\cdots x_{i_n}{:}=H_\nu(\vx)\qquad\text{where $\nu_i$ is the number of times $i$ occurs in $(i_1,\dots,i_n)$}.
  \label{eq:wick-is-hermite}
\end{equation}
The Wick product is symmetric in its indices and depends only on the multiset. Two
degrees are needed explicitly. For degree two, ${:}x_ix_j{:}=x_ix_j-\delta_{ij}$, so that
${:}\vx\vx^{\!\top}{:}=\vx\vx^{\!\top}-\Id$. For degree four, expanding \cref{eq:wick-gf} or
multiplying out $\Hen{k}$'s,
\begin{equation}
\begin{aligned}
  {:}x_ix_jx_kx_l{:}={}&x_ix_jx_kx_l
  -\bigl(\delta_{ij}x_kx_l+\delta_{ik}x_jx_l+\delta_{il}x_jx_k+\delta_{jk}x_ix_l+\delta_{jl}x_ix_k+\delta_{kl}x_ix_j\bigr)\\
  &+\bigl(\delta_{ij}\delta_{kl}+\delta_{ik}\delta_{jl}+\delta_{il}\delta_{jk}\bigr),
\end{aligned}
\label{eq:wick4}
\end{equation}
the inclusion--exclusion over pairs and double pairs; on the diagonal $i=j=k=l$ it
gives $x^4-6x^2+3=\Hen{4}(x)$, as it must. \Cref{eq:wick4} is what the paper means by
``the cumulant contractions with $\delta_{ii}$ are exactly the terms that convert
monomials into Hermite polynomials''.

\begin{lemma}[Second moment of a symmetric contraction]
\label{lem:wick-second-moment}
Let $\mat T$ be a symmetric tensor of order $n$ on $\Reals^{\nsites}$. Then
\begin{equation}
  \E_\gamma\Bigl[\Bigl(\sum_{i_1\cdots i_n}T_{i_1\cdots i_n}\,{:}x_{i_1}\cdots x_{i_n}{:}\Bigr)^2\Bigr]
  =n!\sum_{i_1\cdots i_n}T_{i_1\cdots i_n}^2=n!\,\Frob{\mat T}^2 .
  \label{eq:wick-second-moment}
\end{equation}
\end{lemma}
\begin{proof}
Group the ordered tuples by their multiset $\nu$: each $\nu$ with $|\nu|=n$ collects
$n!/\nu!$ tuples, all with the same entry $T_\nu$ (symmetry) and the same Wick product
$H_\nu$ (\cref{eq:wick-is-hermite}). So the sum is $\sum_{|\nu|=n}(n!/\nu!)T_\nu H_\nu$,
and by \cref{eq:multi-orth} its second moment is
$\sum_\nu(n!/\nu!)^2T_\nu^2\,\nu!=n!\sum_\nu(n!/\nu!)T_\nu^2=n!\sum_{\text{ordered}}T^2$.
\end{proof}

\why{This lemma is where the paper's two constants come from. The leading term of
the score gap has $m$-th component $\tfrac16\sum_{jkl}\kappa_4[mjkl]\,{:}x_jx_kx_l{:}$
times $e^{-4t}$; with $n=3$ and the symmetric tensor $T_{jkl}=\kappa_4[mjkl]$,
\cref{eq:wick-second-moment} gives a $\gamma$-second moment of
$\tfrac1{36}\cdot3!\sum_{jkl}\kappa_4[mjkl]^2$, and summing over $m$ gives
$\tfrac16\Frob{\kfour}^2$. The leading term of a shape score is
$\tfrac1{24}\sum\partial_\theta\kappa_4[ijkl]\,{:}x_ix_jx_kx_l{:}$ times $e^{-4t}$; with
$n=4$, its second moment is $\tfrac{4!}{24^2}\Frob{\partial_\theta\kfour}^2=\tfrac1{24}\Frob{\partial_\theta\kfour}^2$.
The paper's Appendix A states both computations in one line each; this lemma is that
line.}

\section{When is a density ratio in $L^2(\gamma)$, and why Theorem 1 does not care}
\label{sec:l2}

The Hermite expansion \cref{eq:hermite-expansion} of a density ratio $P/\gamma$
requires $P/\gamma\in L^2(\gamma)$, that is
\begin{equation}
  \int\frac{P(\vx)^2}{\gamma(\vx)}\dd\vx<\infty,
  \qquad\text{equivalently}\qquad \chi^2(P\,\|\,\gamma)=\int\frac{P^2}{\gamma}-1<\infty .
\end{equation}
This is a strong condition: it asks $P$ to have tails no heavier than $\sqrt\gamma$,
that is $P(\vx)\lesssim e^{-|\vx|^2/4}$ in every direction.

\begin{proposition}[Gaussian laws]
\label{prop:l2-gaussian}
Let $P=\normal{\vect0}{\Sig}$. Then $P/\gamma\in L^2(\gamma)$ if and only if
$\lmax(\Sig)<2$, and in that case
$\int P^2/\gamma=\prod_k\bigl(\lambda_k(2-\lambda_k)\bigr)^{-1/2}$ over the eigenvalues
$\lambda_k$ of $\Sig$.
\end{proposition}
\begin{proof}
$P^2/\gamma=(2\pi)^{-\nsites/2}(\det\Sig)^{-1}\exp\bigl(-\tfrac12\vx^{\!\top}(2\Sig^{-1}-\Id)\vx\bigr)$.
By \cref{lem:gauss-multi} the integral is finite if and only if $2\Sig^{-1}-\Id\succ0$,
i.e. $\Sig\prec2\Id$, and then equals $(\det\Sig)^{-1}\det(2\Sig^{-1}-\Id)^{-1/2}=\det\bigl(\Sig(2\Id-\Sig)\bigr)^{-1/2}$.
\end{proof}

Two applications, both stated in the paper's Section 3. First, the clean Gaussian
chain at $\corr=0.85$ has $\lmax(\Sig_0)>2$ (the chapter on Markov chains bounds
$\lmax$ between the row sums of the Toeplitz matrix and $(1+\corr)/(1-\corr)$, and
already the $2\times2$ block $\begin{psmallmatrix}1&\corr\\\corr&1\end{psmallmatrix}$ has
eigenvalue $1+\corr>2$ when $\corr>1$, while for $\nsites=32$ the Rayleigh quotient of
the constant vector exceeds $2$ as soon as $\corr\gtrsim0.35$), so $\Pz/\gamma\notin L^2(\gamma)$
for the Gaussian chain itself. Second, the noised law $\Pt=\normal{\vect0}{\Sig_t}$ with
$\Sig_t=e^{-2t}\Sig_0+\Dt\Id$ has eigenvalues $e^{-2t}\lambda_k+1-e^{-2t}$, which are all
below $2$ if and only if $e^{-2t}(\lmax-1)<1$, i.e.
\begin{equation}
  \Pt/\gamma\in L^2(\gamma)\iff e^{2t}>\lmax(\Sig_0)-1 ,
  \label{eq:l2-crossover}
\end{equation}
the exact statement in the paper's paragraph ``Where the asymptote begins''.

\begin{proposition}[The Laplace law]
\label{prop:l2-laplace}
If $P$ is the scalar Laplace density $\tfrac1{2b}e^{-|x|/b}$, then $P/\gamma\notin L^2(\gamma)$
for every $b>0$. The same holds for the Laplace chain on $\Reals^{\nsites}$.
\end{proposition}
\begin{proof}
$P(x)^2/\gamma(x)=\tfrac{\sqrt{2\pi}}{4b^2}\exp\bigl(x^2/2-2|x|/b\bigr)$, and $x^2/2-2|x|/b\to+\infty$,
so the integrand is unbounded and not integrable. For the chain,
$\Pz(\va)=\mu(a_1)\prod_{i\ge2}\innov(a_i-\corr a_{i-1})$ with $\innov$ Laplace; along the
ray $\va=u\vect e_{\nsites}$ (all sites zero except the last) the ratio
$\Pz^2/\gamma$ is proportional to $\exp(u^2/2-2|u|/b)$, and by continuity the same
growth holds on a tube around the ray, so the integral over $\Reals^{\nsites}$ diverges.
\end{proof}

Exponential tails are lighter than Gaussian ones in the sense of moments, and the
Laplace law has an exponential moment; but they are heavier than $e^{-x^2/4}$, and that
is what $L^2(\gamma)$ asks. The same is true of every family in the paper other than
the uniform and the mixture.

\paragraph{Why the paper's Theorem 1 does not need it.}
A tempting route to attenuation is spectral: the Ornstein--Uhlenbeck semigroup
$(\mathcal P_tf)(a)=\E[f(e^{-t}a+\sqrt{\Dt}Z)]$ is self-adjoint on $L^2(\gamma)$ with
eigenfunctions $\Hen{n}$ and eigenvalues $e^{-nt}$, so one could expand $\Pz/\gamma$ in
the chaos basis and let the semigroup act coefficient by coefficient. That route
needs $\Pz/\gamma\in L^2(\gamma)$, which the two propositions above show fails for the
paper's chain, Gaussian or Laplace. Theorem 1 is instead an identity for one
expectation, $\E_{\Pt}[H_\nu(\vx_t)]$, proved by Fubini for each fixed $\nu$; it needs
only $\E|\va|^{|\nu|}<\infty$, so that $H_\nu(\va)$ is integrable. No expansion of a
density is ever summed. This is why the paper can state the theorem for every law
with finite moments, and why its proofs of the asymptotic theorems go through the
damped cumulant-generating function $\Lambda$ rather than through the chaos expansion of $\Pt/\gamma$.

\section{Mehler's formula and the paper's Theorem 1}

\begin{theorem}[Mehler's formula, moment form]
\label{thm:mehler}
Let $Z\sim\normal01$, $a\in\Reals$, $\alpha\in[0,1]$ and $\beta=\sqrt{1-\alpha^2}$. Then
for every $n\ge0$,
\begin{equation}
  \E_Z\bigl[\Hen{n}(\alpha a+\beta Z)\bigr]=\alpha^n\Hen{n}(a).
  \label{eq:mehler}
\end{equation}
\end{theorem}
\begin{proof}
Put $x=\alpha a+\beta Z$ in the generating function \cref{eq:he-gf} and take the
expectation over $Z$:
\[
  \E_Z\bigl[e^{s(\alpha a+\beta Z)-s^2/2}\bigr]
  =e^{s\alpha a-s^2/2}\,\E_Z\bigl[e^{s\beta Z}\bigr]
  =e^{s\alpha a-s^2/2}\,e^{s^2\beta^2/2}
  =e^{s\alpha a-s^2\alpha^2/2},
\]
using \cref{eq:gauss-moments} and $\beta^2=1-\alpha^2$. The right-hand side is the
generating function \cref{eq:he-gf} evaluated at $x=a$ and parameter $\alpha s$:
$e^{(\alpha s)a-(\alpha s)^2/2}=\sum_n\Hen{n}(a)\alpha^ns^n/n!$. The left-hand side is
$\E_Z[\sum_n\Hen{n}(\alpha a+\beta Z)s^n/n!]$, and the expectation may be taken inside
the sum: the series $\sum_n|\Hen{n}(x)||s|^n/n!$ is dominated, for $|s|\le1$, by a
polynomial in $|x|$ times $e^{|x|}$ (from the explicit coefficient formula,
$|\Hen{n}(x)|\le\sum_k\binom n{2k}(2k-1)!!|x|^{n-2k}\le(1+|x|)^n\,n!!$), which is
integrable against the Gaussian law of $Z$, so dominated convergence applies.
Comparing the coefficients of $s^n/n!$ on the two sides gives \cref{eq:mehler}.
\end{proof}

The name refers to Mehler's 1866 kernel identity,
\[
  \sum_{n\ge0}\Hen{n}(x)\Hen{n}(y)\frac{\alpha^n}{n!}
  =\frac{1}{\sqrt{1-\alpha^2}}\exp\Bigl(-\frac{\alpha^2(x^2+y^2)-2\alpha xy}{2(1-\alpha^2)}\Bigr),
\]
which is the transition density of the Ornstein--Uhlenbeck process divided by
$\gamma(y)$; \cref{eq:mehler} is its moment form, and says that $\Hen{n}$ is an
eigenfunction of the channel with eigenvalue $\alpha^n=e^{-nt}$.

\begin{theorem}[Hermite attenuation; the paper's Theorem 1]
\label{thm:hermite-attenuation}
Let $\vx_t=e^{-t}\va+\sqrt{\Dt}\vz$ with $\va\sim\Pz$ and $\vz\sim\normal{\vect0}{\Id}$
independent. For every multi-index $\nu$ with $\E_{\Pz}|\va|^{|\nu|}<\infty$,
\begin{equation}
  \E_{\Pt}[H_\nu(\vx_t)]=e^{-|\nu|t}\,\E_{\Pz}[H_\nu(\va)].
  \label{eq:attenuation}
\end{equation}
\end{theorem}
\begin{proof}
Condition on $\va$. Given $\va$, the coordinates $x_{t,i}=e^{-t}a_i+\sqrt{\Dt}z_i$ are
independent, each of the form $\alpha a_i+\beta z_i$ with $\alpha=e^{-t}$,
$\beta=\sqrt{1-\alpha^2}=\sqrt{\Dt}$. So
\[
  \E\bigl[H_\nu(\vx_t)\,\big|\,\va\bigr]=\prod_i\E_{z_i}\bigl[\Hen{\nu_i}(\alpha a_i+\beta z_i)\bigr]
  =\prod_i\alpha^{\nu_i}\Hen{\nu_i}(a_i)=\alpha^{|\nu|}H_\nu(\va)
\]
by \cref{thm:mehler} applied coordinate by coordinate. Taking the expectation over
$\va\sim\Pz$ gives \cref{eq:attenuation}, provided $H_\nu(\va)$ is integrable, which
holds because $|H_\nu(\va)|\le C(1+|\va|)^{|\nu|}$ and $\E|\va|^{|\nu|}<\infty$; the
tower property $\E[\E[\cdot\mid\va]]=\E[\cdot]$ is Fubini's theorem, which also needs
only that integrability.
\end{proof}

The theorem is exact at every $t$, needs no assumption on the dependence between
sites, and needs no $L^2(\gamma)$ condition. It is verified by Monte Carlo in
\code{tests/test\_theory\_hermite.py} on a Laplace sample, and the cases $n\le4$ are
checked in Lean as polynomial identities under the first four Gaussian moments
(\code{lean/README.md}).

\section{Two exact consequences}

\begin{corollary}[Cumulants of order three and higher]
\label{cor:cumulant-attenuation}
For every $m\ge3$ with $\E|\va|^m<\infty$, the order-$m$ cumulant tensor of $\vx_t$ is
$e^{-mt}$ times that of $\va$. The covariance is $\Cov(\vx_t)=e^{-2t}\Sig_0+\Dt\Id$.
\end{corollary}
\begin{proof}
$\vx_t=e^{-t}\va+\sqrt{\Dt}\vz$ is a sum of two independent vectors. By additivity and
the transformation rule of the chapter on moments and cumulants, the order-$m$
cumulant tensor of $\vx_t$ is $e^{-mt}$ times that of $\va$ plus $\Dt^{m/2}$ times that
of $\vz$; the Gaussian $\vz$ has no cumulant tensor of order $m\ge3$. For $m=2$ the
Gaussian contributes $\Dt\Id$.
\end{proof}

This corollary uses no Hermite polynomials, only the arithmetic of cumulants; the
paper's parenthesis ``Gaussian noise has none'' is the whole proof. The Hermite form
of the theorem is nevertheless the one the paper needs, because Fisher information is
built from moments, and the next corollary says how the Hermite moments of degree
four see the fourth cumulant.

\begin{corollary}[Degree-four Hermite moments and the fourth-cumulant tensor]
\label{cor:degree-four}
Let $\va$ have mean zero, covariance $\Sig_0$ and finite fourth moments. For any
indices, with $\nu$ the multiset of $(i,j,k,l)$,
\begin{equation}
  \E[H_\nu(\va)]=\E[{:}a_ia_ja_ka_l{:}]
  =\kappa_4[ijkl]+(\Sig_0-\Id)_{ij}(\Sig_0-\Id)_{kl}+(\Sig_0-\Id)_{ik}(\Sig_0-\Id)_{jl}+(\Sig_0-\Id)_{il}(\Sig_0-\Id)_{jk}.
  \label{eq:he4-kappa4}
\end{equation}
Consequently, if two laws have the same covariance $\Sig_0$, their degree-four
Hermite moments differ exactly by the difference of their fourth-cumulant tensors,
and after noising that difference is $e^{-4t}(\kfour-\kfour')$.
\end{corollary}
\begin{proof}
Take the expectation of \cref{eq:wick4}. By the fourth-cumulant formula of the chapter
on moments and cumulants, $\E[a_ia_ja_ka_l]=\kappa_4[ijkl]+\Sigma_{ij}\Sigma_{kl}+\Sigma_{ik}\Sigma_{jl}+\Sigma_{il}\Sigma_{jk}$
with $\Sigma=\Sig_0$; the six pair terms give $-\delta_{ij}\Sigma_{kl}-\dots$; the three
double-pair terms give $+\delta_{ij}\delta_{kl}+\dots$. Collecting by pairing,
$\Sigma_{ij}\Sigma_{kl}-\delta_{ij}\Sigma_{kl}-\Sigma_{ij}\delta_{kl}+\delta_{ij}\delta_{kl}=(\Sigma-\Id)_{ij}(\Sigma-\Id)_{kl}$,
and likewise for the other two pairings. The second statement follows because the
product terms depend only on $\Sig_0$, and the attenuation by $e^{-4t}$ is
\cref{thm:hermite-attenuation} with $|\nu|=4$.
\end{proof}

For the paper's variance-matched families, $\Sig_0$ is the same for all, so the
degree-four Hermite moment of the chain differs from that of its covariance twin by
exactly $\kfour$; on the diagonal, $\E[\Hen{4}(a_i)]=\kappa_4(a_i)$ because $\Var(a_i)=1$.
Degree two works the same way: $\E[{:}a_ia_j{:}]=(\Sig_0-\Id)_{ij}$, attenuated by $e^{-2t}$,
which is the ``correlation coordinate'' and is consistent with
\cref{cor:cumulant-attenuation}. This is the paper's second exact consequence, checked
in Lean as target 2 (\code{lean/README.md}).

\section{The Edgeworth expansion: the first correction}
\label{sec:edgeworth}

The paper's E16 asks how much of the gap between the true score and the twin's is
explained by ``adding the fourth-order term back''. The object added is the first
term of the Edgeworth (Gram--Charlier A) expansion, defined here; the chapter on
misspecification and Edgeworth corrections uses it.

\paragraph{Scalar form.}
Let $X$ have mean $0$, variance $1$ and cumulants $\kappa_3,\kappa_4,\dots$. Formally,
its generating function is $\exp(\tfrac12s^2+\sum_{n\ge3}\kappa_ns^n/n!)$, that of the
standard Gaussian times $\exp(\sum_{n\ge3}\kappa_ns^n/n!)$; since multiplying a
generating function by $s^n$ corresponds to applying $(-\dd/\dd x)^n$ to the density,
and $(-\dd/\dd x)^n\gamma=\Hen{n}\gamma$ by the Rodrigues formula \cref{eq:rodrigues},
\begin{equation}
  p(x)\ \text{``$=$''}\ \exp\Bigl(\sum_{n\ge3}\frac{\kappa_n}{n!}\bigl(-\tfrac{\dd}{\dd x}\bigr)^n\Bigr)\gamma(x)
  =\gamma(x)\Bigl[1+\frac{\kappa_3}{6}\Hen{3}(x)+\frac{\kappa_4}{24}\Hen{4}(x)+\dots\Bigr],
  \label{eq:edgeworth-scalar}
\end{equation}
where the dots collect $\kappa_3^2\Hen{6}/72$ and all higher terms. The
\term{first Edgeworth correction} is the truncation
$p\approx\gamma(1+\kappa_3\Hen{3}/6+\kappa_4\Hen{4}/24)$, and for a symmetric law it
is $p\approx\gamma(1+\kappa_4\Hen{4}/24)$. The series need not converge (it does not
for the Laplace law, whose $\kappa_n$ grow like $n!$), and even when it does the
truncation can be negative for large $|x|$; the expansion is an asymptotic device,
exact order by order in the cumulants, and the paper reports where it helps and where
it diverges rather than assuming either.

\paragraph{Multivariate form around a Gaussian with covariance $\Sig$.}
Replace $\gamma$ by $\gamma_{\Sig}$, the density of $\normal{\vect0}{\Sig}$, and the
Hermite polynomials by the $\Sig$-\term{Hermite tensors}
\[
  h_{i_1\cdots i_n}(\vx)=(-1)^n\,\frac{\partial_{i_1}\cdots\partial_{i_n}\gamma_\Sig(\vx)}{\gamma_\Sig(\vx)} .
\]
With $\vect y=\Sig^{-1}\vx$, $\partial_i\gamma_\Sig=-y_i\gamma_\Sig$ and
$\partial_iy_j=(\Sig^{-1})_{ij}$, so the fourth tensor is
\begin{equation}
\begin{aligned}
  h_{ijkl}(\vx)={}&y_iy_jy_ky_l-\bigl[(\Sig^{-1})_{ij}y_ky_l+\text{five more}\bigr]\\
  &+\bigl[(\Sig^{-1})_{ij}(\Sig^{-1})_{kl}+(\Sig^{-1})_{ik}(\Sig^{-1})_{jl}+(\Sig^{-1})_{il}(\Sig^{-1})_{jk}\bigr],
\end{aligned}
\end{equation}
\cref{eq:wick4} with $\vx\mapsto\vect y$ and $\delta\mapsto\Sig^{-1}$. For a symmetric
law with fourth-cumulant tensor $\mat K$ the first correction is
$p\approx\gamma_\Sig(1+\epsilon)$ with $\epsilon=\tfrac1{24}\sum_{ijkl}K_{ijkl}h_{ijkl}$;
contracting a symmetric $\mat K$ with the three groups of terms gives, with
$A_{kl}=\sum_{ij}K_{ijkl}(\Sig^{-1})_{ij}$,
\begin{equation}
  \epsilon(\vx)=\frac1{24}\Bigl[\mat K[\vect y,\vect y,\vect y,\vect y]-6\,\vect y^{\!\top}\mat A\vect y+3\langle\mat A,\Sig^{-1}\rangle\Bigr],
  \qquad
  \nabla\epsilon=\frac16\Sig^{-1}\mat K[\cdot,\vect y,\vect y,\vect y]-\frac12\Sig^{-1}\mat A\vect y,
  \label{eq:edgeworth-multi}
\end{equation}
where $\mat K[\vect y,\vect y,\vect y,\vect y]=\sum K_{ijkl}y_iy_jy_ky_l$ and
$\langle\mat A,\Sig^{-1}\rangle=\sum_{kl}A_{kl}(\Sig^{-1})_{kl}$; the six pair terms
each contract to $\vect y^{\!\top}\mat A\vect y$ and the three double-pair terms each to
$\langle\mat A,\Sig^{-1}\rangle$ by the symmetry of $\mat K$, and the gradient uses
$\nabla_{\vx}=\Sig^{-1}\nabla_{\vect y}$. The score of the corrected density is
$\nabla\log[\gamma_\Sig(1+\epsilon)]=-\Sig^{-1}\vx+\nabla\epsilon/(1+\epsilon)$. This is
exactly the formula in the paper's Appendix E for E16, with $\Sig=\Sig_t$ and
$\mat K=e^{-4t}\kfour$: by \cref{cor:cumulant-attenuation} the noised chain has
covariance $\Sig_t$ and fourth-cumulant tensor $e^{-4t}\kfour$, so the first
Edgeworth correction around the covariance twin is the degree-four Hermite term that
the paper's Theorem 2 identifies as the whole of the leading gap.

\section{Worked example: a unit-variance Laplace site at $t=\log2$}

Let $a$ have the unit-variance Laplace law, $b=1/\sqrt2$, with $\kappa_4=3$ (chapter on
moments and cumulants), and let $x_t=e^{-t}a+\sqrt{\Dt}z$ at $t=\log2$, so
$\alpha=e^{-t}=\tfrac12$ and $\Dt=\tfrac34$. Theorem 1 predicts
\[
  \E[\Hen{4}(x_t)]=\alpha^4\E[\Hen{4}(a)]=\tfrac1{16}\bigl(\E a^4-6\E a^2+3\bigr)=\tfrac1{16}(6-6+3)=\tfrac3{16},
\]
using $\E a^4=24b^4=6$. Directly: $\E x_t^2=\tfrac14+\tfrac34=1$ and
$\E x_t^4=\alpha^4\E a^4+6\alpha^2\Dt\E a^2\E z^2+\Dt^2\E z^4=\tfrac6{16}+6\cdot\tfrac14\cdot\tfrac34+\tfrac9{16}\cdot3=\tfrac{6+18+27}{16}=\tfrac{51}{16}$,
so $\E[\Hen{4}(x_t)]=\tfrac{51}{16}-6+3=\tfrac3{16}$, as predicted. The fourth cumulant of
$x_t$ is $\kappa_4(x_t)=\E x_t^4-3=\tfrac3{16}=e^{-4t}\kappa_4(a)$, illustrating
\cref{cor:cumulant-attenuation}; and the degree-two moment $\E[\Hen{2}(x_t)]=\E x_t^2-1=0$
equals $e^{-2t}\E[\Hen{2}(a)]=0$, as unit variance requires.

For the constant of \cref{lem:wick-second-moment} in the scalar case: the leading term
of the score gap is $(\kappa_4/6)\Hen{3}(x)e^{-4t}$, whose $\gamma$-second moment is
$(\kappa_4/6)^2\cdot3!\,e^{-8t}=(\kappa_4^2/6)e^{-8t}$, that is $\tfrac96e^{-8t}=1.5\,e^{-8t}$
for the Laplace law; this is the value that \code{tests/test\_theory\_hermite.py}
compares with a fine quadrature of the scalar noised density.

\where{The paper's Section 3 (Theorem 1 and its two consequences, the Hermite
normalisation $\nu!$, the $L^2(\gamma)$ remarks, the Edgeworth paragraph ``The gap is
the fourth-order term''), Appendix A (the proofs of Theorems 1 and 2, Wick's theorem
for the constants, the crossover $e^{2t}>\lmax-1$) and Appendix E (E16). Experiments
E01, E16, E18, E21. Code: the functions \code{he}, \code{hermite\_moment\_attenuation},
\code{score\_gap\_constant} and \code{shape\_information\_constant} of
\code{src/markov\_score/theory.py}; the Edgeworth oracle
\code{experiments/e16\_edgeworth\_oracle.py}; the tests in
\code{tests/test\_theory\_hermite.py}; and the Lean targets listed in
\code{lean/README.md}.}

\paragraph{References.}
G. Szeg\H{o} (1975), \emph{Orthogonal Polynomials}, 4th edition, American Mathematical
Society (Chapter V: Hermite polynomials, generating function, Mehler's formula).
L. Isserlis (1918), ``On a formula for the product-moment coefficient of any order of a
normal frequency distribution in any number of variables'', Biometrika 12, 134--139.
F. G. Mehler (1866), ``Ueber die Entwicklung einer Function von beliebig vielen
Variablen nach Laplaceschen Functionen h\"oherer Ordnung'', Journal f\"ur die reine und
angewandte Mathematik 66, 161--176.
S. Janson (1997), \emph{Gaussian Hilbert Spaces}, Cambridge University Press (Wick
products, chaos expansion, the Ornstein--Uhlenbeck semigroup on $L^2(\gamma)$).
P. Hall (1992), \emph{The Bootstrap and Edgeworth Expansion}, Springer (Chapter 2:
the Edgeworth expansion and its first terms).
P. McCullagh (1987), \emph{Tensor Methods in Statistics}, Chapman and Hall (Chapter 5:
Edgeworth series in several variables, Hermite tensors).
